\cotefile{137}
\scanpage{2}
% archive foliation: 1
\margin{\unsure{le même} \illegible{3}; \unsure{$k_*$} \illegible{5}}

\uline{Lemme 1} Soit $L_*$ un $k$-module $\frac12$ \unsure{simplicial}.
Pour que $L_*$ soit un \unsure{$k_*$-module} projectif, il suffit
qu'il soit flasque (i.e. $\unsure{\pi_i}(L_*)=0$ $\forall i\geq 0$), et que
les $L_n$ soient des $k$-modules projectifs.
\margin{\unsure{NB} \illegible{2} \unsure{$L_*$} \illegible{3} \unsure{projectif}}

La \unsure{cond°} \illegible{1} \unsure{signifie aussi} que le complexe de
chaînes (de $k$-modules) associé par Dold-Puppe
est acyclique et à \unsure{composantes projectives}. On voit
alors qu'il est \inserted{\unsure{donc}} isomorphe à une somme directe
\unsure{de complexes de chaînes de la forme}
\[
\unsure{0\longrightarrow L\xrightarrow{\mathrm{id}}L\longrightarrow 0}.
\]
\unsure{On} \illegible{3} \unsure{tel}
\illegible{2} \unsure{correspond} \illegible{4}
\unsure{complexe} \illegible{2} \unsure{projectif} \illegible{3}
\unsure{au} \illegible{3} (\illegible{3})
\illegible{5} \unsure{donc on} \illegible{3}.

\uline{NB} Il suffit \unsure{en fait} \illegible{2} \unsure{que}
$L_*$ \unsure{vérifie} i.e. $\pi_i(L_*)=0$ $\forall i>0$, \unsure{au lieu de}
$\pi_i(L_*)=0$ pour $i\geq0$, et $\pi_0(L_*)$ $k$-projectif.
(\unsure{On trouve les mêmes conclusions} \illegible{2}
\unsure{c. et s.}) En effet, dans l'argument précédent,
il faut de plus tenir compte du \unsure{som-}
\unsure{mande} qui est un module projectif
$L_0$ placé au degré $0$. Or il est \unsure{immédiat}
que c'est aussi un complexe de \unsure{$k_*$-}
chaînes projectif.

\uline{Cor.} Les $\unsure{\Gamma^i\Phi_{* k}}\otimes_k\unsure{\Lambda^j\Psi_{* k}}$
\inserted{et $\unsure{\Lambda^i\Phi_{* k}}\otimes\unsure{\Lambda^j\Psi_{* k}}$},
(pour $i>0$) sont des
modules \unsure{simpliciaux} projectifs.
